\subsection{The special Cohen weight does not lower the heat-cycle filtration}

The following calculation addresses the specific automatic family
$b=(p+3)/2$.  It does not rule out a stronger first-coefficient
theorem obtained by other arithmetic methods.  It does rule out
lowering the current Jacobi heat construction to bounded weight
by Hasse division or by moving around its heat cycle.

\begin{proposition}[KT-COHEN-SPECIAL-HEAT-FILTRATION]
Let $p\geq31$ be prime with $p\equiv3\pmod4$, put
$b=(p+3)/2$ and $\kappa=b+1=(p+5)/2$, and let
\[
 \phi=\zeta(1-2b)E_{\kappa,1}
\]
be the normalized Jacobi Eisenstein series used in
KT-JACOBI-TWO-UNIT-FIELDS.  Let $L=L_1$ be the heat operator,
which multiplies the coefficient at $q^n\zeta^r$ by $4n-r^2$.
Write $\Omega$ for Jacobi filtration modulo $p$ at index $1$.
Then
\[
 \Omega(\phi)=\kappa,\qquad
 \Omega(L^i\phi)=\kappa+i(p+1)\quad(1\leq i\leq p-2),
 \tag{HC.1}
\]
while $L^{p-1}\phi\ne0$ modulo $p$ and
\[
 \Omega(L^{p-1}\phi)\geq\kappa.
 \tag{HC.2}
\]
For $i\geq1$ the reductions of $L^i\phi$ are periodic in $i$
with period $p-1$.  Thus every nonzero iterate has filtration
at least $\kappa$, and the first heat lift has exact filtration
\[
 \Omega(L\phi)=\frac{3p+7}{2}.
 \tag{HC.3}
\]
\end{proposition}

\begin{proof}
The $p$-integrality of $\phi$ is part of the already checked
Jacobi construction.  Its constant is
\[
 \zeta(1-2b)=-\frac{B_{p+3}}{p+3}
             \equiv-\frac{B_4}{4}=\frac1{120}\pmod p,
 \tag{HC.4}
\]
by Kummer's congruence.  The Euler factors at $p$ in this
congruence are units with residue $1$.  Hence $\phi$ is nonzero.

The congruent-weight theorem for Jacobi forms of fixed index
implies that every weight representing its nonzero reduction
is congruent to $\kappa$ modulo $p-1$.  Since
$0<\kappa<p-1$ and Jacobi forms of negative weight vanish,
this proves $\Omega(\phi)=\kappa$.

For a nonzero reduction of index $1$, the heat-filtration
theorem states
\[
 \Omega(Lf)\leq\Omega(f)+p+1,
 \quad\hbox{with equality precisely when }
 p\nmid 2\Omega(f)-1.
 \tag{HC.5}
\]
It also supplies a Jacobi modular lift of the heat image.
For $0\leq i\leq p-3$, if the claimed value holds at $i$, then
\[
 2\bigl(\kappa+i(p+1)\bigr)-1
       \equiv4+2i\not\equiv0\pmod p.
 \tag{HC.6}
\]
The unique first index in $0,\ldots,p-1$ where this expression
vanishes is $i=p-2$.  Induction by (HC.5) therefore proves
(HC.1), including nonvanishing at every step.

Coefficientwise Fermat gives $L^p\phi=L\phi$ modulo $p$.
In particular $L^{p-1}\phi$ cannot vanish, since $L\phi$
does not vanish.  Apply the inequality in (HC.5) to that
nonzero reduction:
\[
 \kappa+p+1=\Omega(L^p\phi)
       \leq\Omega(L^{p-1}\phi)+p+1.
 \]
Subtracting $p+1$ proves (HC.2).
Finally $L^{i+p-1}\phi=L^i\phi$ for every $i\geq1$, again
coefficientwise.  This proves all the assertions.
\end{proof}

\paragraph{Meaning for the short-discriminant proposal.}
The first heat image is the one used to remove the
discriminant-zero coefficients before applying the odd-component
and Sturm argument.  Equation (HC.3) shows that its linear
weight cannot be lowered by dividing out a Hasse invariant
while remaining a holomorphic Jacobi form with the same
reduction.  Iterating the heat operator never reaches a
filtration smaller than $(p+5)/2$ either.  This excludes that
specific way of replacing the existing $O(p)$ bound by the
$O(p^{2/5})$ sufficient window in KT-BETA3-SHORT-WINDOW.
It does not give a lower bound on the actual first nonzero
fundamental coefficient.

The nonzero constant term alone is not a field coefficient.
In Jacobi notation, its translates have $4n-r^2=0$; the heat
operator kills all of them.  The proof needs a nonzero
coefficient with positive discriminant and with odd $r$.
Thus the constant in (HC.4) cannot simply be counted as the
desired seed.

The fixed-weight identity
$\beta_b(\chi_D)\equiv\beta_2(\chi_D\chi_{-p})$
does not contradict (HC.1).  Its real quadratic character has
conductor $p|D|$, and the corresponding weight-$5/2$ Cohen
coefficient is at index $p|D|$, not at index $|D|$.
A small ordinary Sturm bound for the fixed-weight form
does not bound its first unit coefficient in this
$p$-dependent subsequence.
For the short window itself the class-number requirement
would present no extra gap: $d^5\leq1600p^2$ implies
$d<4p$ for $p\geq31$, and the carried
KT-SMALL-DISCRIMINANT-CLASS-BOUND gives $h(-d)<p$.
Existence of a $\beta_b$-unit field in that window remains
unproved here.

\paragraph{Exact primary input and review scope.}
D.~Choi, Y.~Choie and O.~K.~Richter,
\emph{Congruences for Siegel modular forms of degree two},
Ann.~Inst.~Fourier \textbf{61} (2011), 1455--1466:
Section 2, Propositions 1 and 2 on printed p.~1457 give
the congruent-weight theorem and (HC.5).
Their hypotheses are even Jacobi weight, $p\geq5$, and
fixed index; here $\kappa$ and all lift weights are even
and the index is $1$.  The definitions and both full
statements, with the surrounding scope, were read directly
in the retained primary text
\texttt{main\_gate/literature/character\_unit\_followup/Choi\_Choie\_Richter\_2011}.
These published propositions are imported inputs; their
foundational Sofer and Richter proofs are not newly
reconstructed in this deduction.  The boundary case
$p=2\kappa-5$ is not inserted into either strict inequality
in their Theorem 2.
